\begin{textsample}{POS Dim 1 – es – Score 40.00 – es/es\_em\_4.txt}  \label{ex:f1_pos_198}
A Formação do Novo Ensino Médio O documento curricular da etapa do Ensino Médio se constituirá de duas partes : uma relacionada à formação geral básica, que estabelece um conjunto de \textbf{competências} e habilidades das áreas de conhecimento previstas na BNCC, e outra parte relacionada aos itinerários formativos : conjunto de unidades curriculares ofertado pelas instituições e redes de ensino que possibilitam ao estudante aprofundar seus conhecimentos e se preparar para o prosseguimento de estudos ou para o mundo do trabalho, de forma a contribuir para a construção de soluções de \textbf{problemas} específicos da sociedade ( BRASIL, 2018a ).
Na formação do Novo Ensino Médio encontram -se a pesquisa como prática pedagógica ; a \textbf{compreensão} da diversidade e realidades dos sujeitos ; a diversificação da oferta de forma a possibilitar múltiplas trajetórias por parte dos estudantes, além da dissociabilidade entre educação e prática no processo de \textbf{ensino-aprendizagem}.
Diante desses fatos, o Ensino Médio prescrito da formação geral básica desse Currículo requer um caráter mediador de todo processo de transformação, por meio do estímulo ao protagonismo estudantil, promoção de atividades práticas e metodologias que fujam de \textbf{meras} \textbf{memorizações}, uma vez que as informações já se encontram à disposição dos sujeitos em tempo real.
Nesse sentido, é \textbf{indispensável} que sejam desenvolvidas, nos estudantes, \textbf{capacidades} de abstração, reflexão, interpretação, proposição e ação, possibilitando a \textbf{compreensão} dos fenômenos naturais e culturais, estimulando a sua autonomia e \textbf{capacidade} de tomar decisões fundamentadas e responsáveis, assim como a promoção de atitudes cooperativas e propositivas para o enfrentamento dos desafios de sua comunidade, do mundo do trabalho e da sociedade em geral.
É importante ressaltar que o Currículo do Espírito Santo, na formação geral básica, considera a formação do sujeito de forma significativa, substituindo a concepção de acúmulo de conceitos e fatos pela \textbf{compreensão} de aprendizagens essenciais, conforme descrita nas Diretrizes Curriculares Nacionais – DCNs ( BRASIL, 2018b ) : > Art. 7º [... ] § 3º As aprendizagens essenciais são as que desenvolvem \textbf{competências} e habilidades entendidas como conhecimento em ação, com significado para a vida, expressas em práticas cognitivas, profissionais e \textbf{socioemocionais}, atitudes e valores continuamente \textbf{mobilizados}, articulados e integrados, para resolver \textbf{demandas} \textbf{complexas} da vida cotidiana, do exercício da cidadania e da atuação no mundo do trabalho.
Essa reestruturação curricular interfere também na articulação dos saberes das diversas áreas de conhecimento, em consonância com o que foi estabelecido pelas DCNs.
A integração curricular, na perspectiva da formação integral do estudante, foi considerada pelo Currículo do Espírito Santo, no qual redatores das diversas áreas e componentes curriculares se dedicaram para construir um documento que dialogasse com os saberes, de \textbf{maneira} a estabelecer o conhecimento integrado dentro de cada área.
Por meio de diferentes pontos de encontros, os conhecimentos específicos de cada componente se cruzam e se correlacionam, de modo a construir um novo e mais \textbf{complexo} conhecimento.
Outra função importante do Ensino Médio, que muitas vezes é atribuída exclusivamente à família, é propiciar o aprimoramento do educando como pessoa humana, dotada de valores como empatia, solidariedade, voltado a práticas não violentas ou discriminatórias, viabilizando, desse modo, que se torne um cidadão atuante na construção de uma sociedade mais justa, democrática e inclusiva.
Nesse ponto, pode ser observado um diálogo com as ideias de Morin ( 2013 ) sobre as necessidades para a educação futura envolvendo diferentes saberes.
Segundo o \textbf{autor}, deve -se fornecer aos estudantes uma cultura que lhes permita articular, religar e contextualizar, situando tais estudantes em um contexto que oportunize reunir os conhecimentos que já foram adquiridos, não sendo \textbf{pertinente} um conhecimento totalmente isolado. \textbf{Tanto} na BNCC quanto no Currículo do Espírito Santo são tratados alguns desses saberes, a exemplo do conhecimento \textbf{pertinente} que se refere ao conhecimento contextualizado, sendo utilizado como base o pensamento \textbf{complexo} e não a visão fragmentada do todo, uma vez que quanto mais fragmentada é a realidade, menor é a \textbf{capacidade} humana de buscar soluções para \textbf{problemas} mais amplos, visto que estão limitados a conhecimentos de apenas parte da realidade no qual estão inseridos.
Uma vez reconhecido esse papel da educação, é necessário que as escolas presentes no território capixaba reverenciem, à luz da BNCC ( BRASIL, 2018a ), algumas atitudes, tais como : sejam espaços que permitam aos estudantes valorizar ações que estimulem o diálogo e não a violência ; o respeito à dignidade do outro, favorecendo o convívio entre os diferentes, sem que haja qualquer alusão à ideia de que alguns grupos tenham maior valor, ou sejam menos dignos \textbf{devido} a sua diferença ; o combate à discriminação e às violações a pessoas ou grupos sociais, visto que na construção cultural histórica do Espírito Santo temos diferentes representações, como Pomeranos, Quilombolas, Indígenas, entre outros grupos que constituem a identidade do estado ; o respeito e a valorização de diferentes manifestações religiosas, a exemplo das religiões de \textbf{matriz} africana, como umbanda e candomblé, das quais muitos estudantes fazem parte e não são contemplados em nenhum momento das atividades escolares, como as células de \textbf{alunos} protestantes e algumas atividades da \textbf{juventude} católica ; a participação política e social, visto ser o jovem parte fundamental do processo de transformação política e social ; e a construção de projetos pessoais e coletivos, \textbf{baseados} na liberdade, na justiça social, na solidariedade e na sustentabilidade.
Tendo em \textbf{vista} as recomendações relacionadas à diversidade, \textbf{sugeridas} pelo Conselho Nacional de Educação, em seu parecer CNE/CP N{\EmojiFont °} 11/2009, onde é estimulada a construção de currículos flexíveis, devem ser propostos itinerários formativos diversificados, que se atentem à heterogeneidade e à \textbf{pluralidade} dos jovens e respondam a seus interesses e necessidades.
Desse modo, o Currículo do Espírito Santo \textbf{segue} a mesma direção da BNCC, norteada pela Lei nº 13.415/2017 que alterou a LDB, estabelecendo que : > Art. 36- O currículo do ensino do ensino \textbf{médio} será composto pela Base Nacional Comum Curricular e por itinerários formativos, que deverão ser organizados por meio da oferta de diferentes arranjos curriculares, conforme a relevância para o contexto local e a possibilidade dos sistemas de ensino, a saber : > > I- linguagens e suas \textbf{tecnologias} ; > II- matemática e suas \textbf{tecnologias} ; > III- \textbf{ciências} da natureza e suas \textbf{tecnologias} ; > IV- \textbf{ciências} humanas e sociais aplicadas ; > V- formação técnica e profissional.
Outro aspecto importante a ser abordado é a preparação básica para o mundo do trabalho, o que não significa retornar ao ensino tecnicista proposto nos anos 70, que representou o reflexo da realidade da \textbf{sistematização} político-social sobre a educação ( SAVIANI, 2008 ), mas sim fazer da escola o ambiente no qual possam ser desenvolvidas \textbf{competências} e habilidades que \textbf{auxiliem} na inserção do estudante de forma ativa, \textbf{crítica}, criativa e responsável em mundo do trabalho cada vez mais seletivo, exigente, \textbf{complexo} e imprevisível.
De acordo com as DCNs, atualizada pela resolução nº 3, de 21 de novembro de 2018, o trabalho é conceituado na sua perspectiva ontológica de transformação da natureza, ampliada como impulsionador do desenvolvimento cognitivo, como realização \textbf{inerente} ao ser humano e como mediação no processo de produção da sua existência ( BRASIL, 2018b ).
Para \textbf{tanto}, é necessário que os projetos pedagógicos e os currículos, entre eles o do Espírito Santo, estejam estruturados, assim como contextualiza na BNCC ( 2018 ), de modo a : \textbf{explicitar} que o trabalho produz e transforma a cultura e modifica a natureza ; propiciar o desenvolvimento de atividades práticas, aliadas à fundamentação \textbf{teórica} em diferentes áreas do conhecimento, possibilitando que o estudante possa usar esse conhecimento em seu Projeto de Vida e promover o \textbf{aprofundamento} em campos de interesse dos \textbf{discentes}, para que as \textbf{competências} e habilidades desenvolvidas no ambiente escolar possam ter algum valor a ser \textbf{agregado} com experiência profissional.
Visto que o mundo do trabalho \textbf{exige} experiência e qualificação, deve -se associar as atividades propostas no ambiente escolar com a realidade do \textbf{mercado} de trabalho capixaba, a exemplo das atividades pesqueiras, agropecuária, entre outras áreas voltadas para a utilização das novas \textbf{tecnologias}, \textbf{explicitar} que a preparação para o mundo do trabalho não está diretamente ligada à profissionalização precoce dos jovens – uma vez que eles viverão em um mundo com profissões e ocupações hoje desconhecidas, caracterizadas pelo uso intensivo de \textbf{tecnologias} –, mas à abertura de possibilidades de atuação \textbf{imediata}, a \textbf{médio} e a longo prazos e para solução de novos \textbf{problemas}.
A flexibilização dos itinerários também favorece o protagonismo, visto que serão ofertados variados itinerários formativos para atender a multiplicidade e interesse dos estudantes, além de permitir ao estado a autonomia necessária para as adequações do currículo, de modo a respeitar as singularidades presentes no território capixaba, sendo expresso nos itinerários formativos, nas unidades curriculares eletivas e/ou grupos de estudo e no desenvolvimento do Projeto de Vida, visando a atender as particularidades dos estudantes e das regiões nas quais estão inseridos.

% matched lemmas: agregar, aluno, aprofundamento, autor, auxiliar, basear, capacidade, ciência, competência, complexo, compreensão, crítico, demanda, devido, discente, ensino-aprendizagem, exigir, explicitar, imediato, indispensável, inerente, juventude, maneira, matriz, memorização, mercado, mero, mobilizar, médio, pertinente, pluralidade, problema, segar, sistematização, socioemocional, sugerir, tanto, tecnologia, teórico, vista
\end{textsample}
